\documentclass[11pt]{article}

\usepackage{amsmath, amssymb, amsfonts}
\usepackage{geometry}
\geometry{margin=1in}

\title{Supplementary Information: Information Distortion, Effective Sample Size, Corrected Covariance, and Bias}
\author{}
\date{}

\begin{document}

\maketitle

\section*{1. Log-Likelihood and Score Definitions}

Consider a trajectory of $T$ sequential measurement intervals.
The total log-likelihood is defined as

\begin{equation}
\ell(\theta) = \log L(\theta) = \sum_{t=1}^T \ell_t(\theta),
\end{equation}

with total score

\begin{equation}
S(\theta) = \nabla_\theta \ell(\theta)
= \sum_{t=1}^T s_t(\theta),
\qquad
s_t(\theta) = \nabla_\theta \ell_t(\theta).
\end{equation}

Expectations and covariances are taken over independent realizations of full trajectories generated under parameter $\theta$.

\section*{2. Gaussian Fisher Information (GFI)}

At each interval, the predictive likelihood is approximated as Gaussian:

\begin{equation}
y_t \sim \mathcal{N}\!\left(\mu_t(\theta), \sigma_t^2(\theta)\right),
\end{equation}

where $\mu_t(\theta)$ and $\sigma_t^2(\theta)$ are obtained recursively.

For scalar variance $\sigma_t^2$, the per-interval Gaussian Fisher information is

\begin{equation}
\mathbf I_t^{\mathrm{G}}(\theta)
=
\frac{1}{\sigma_t^2(\theta)} 
\nabla_\theta \mu_t(\theta)
\nabla_\theta \mu_t(\theta)^\top
+
\frac{1}{2\sigma_t^4(\theta)} 
\nabla_\theta \sigma_t^2(\theta)
\nabla_\theta \sigma_t^2(\theta)^\top.
\end{equation}

The total Gaussian Fisher Information is

\begin{equation}
\mathbf H(\theta)
=
\sum_{t=1}^T \mathbb{E}\!\left[\mathbf I_t^{\mathrm{G}}(\theta)\right]
=
- \mathbb{E}\!\left[\nabla_\theta^2 \ell(\theta)\right].
\end{equation}

\section*{3. Score Covariance Matrix (SCM)}

\begin{equation}
\mathbf J(\theta)
=
\mathrm{Cov}\!\left(S(\theta)\right).
\end{equation}

This matrix includes both per-interval score variance and cross-interval covariances.

\section*{4. Information Distortion Matrix (IDM)}

\begin{equation}
\mathbf C(\theta)
=
\mathbf H^{-1/2}
\mathbf J
\mathbf H^{-1/2}.
\end{equation}

If the predictive model is fully consistent:

\begin{equation}
\mathbf J = \mathbf H,
\qquad
\mathbf C = \mathbf I.
\end{equation}

\section*{5. Per-Parameter Distortion}

The diagonal entries $C_{ii}$ quantify parameter-wise variance inflation relative to the Gaussian Fisher prediction.

\section*{6. Distortion-Corrected Covariance (DCC)}

\begin{equation}
\Sigma_{\mathrm{DCC}}
=
\mathbf H^{-1}
\mathbf J
\mathbf H^{-1}
=
\mathbf H^{-1/2}
\mathbf C
\mathbf H^{-1/2}.
\end{equation}

\section*{7. Distortion-Induced Bias (DIB)}

If

\begin{equation}
\mathbf g = \mathbb{E}[S(\theta)] \neq 0,
\end{equation}

then

\begin{equation}
\mathbf b_{\mathrm{DIB}}
=
\mathbf H^{-1} \mathbf g.
\end{equation}

\section*{8. Effective Sample Size and Residual Dependence}

\subsection*{8.1 Variance of Sum vs Sum of Variances}

For a sequence $X_t$ define:

\begin{equation}
L_X = \sum_{t=1}^T X_t.
\end{equation}

Then:

\begin{equation}
\mathrm{Var}(L_X)
=
\sum_{t=1}^T \mathrm{Var}(X_t)
+
2\sum_{i<j}\mathrm{Cov}(X_i,X_j).
\end{equation}

Define the inflation factor:

\begin{equation}
\kappa_X
=
\frac{\mathrm{Var}(L_X)}
{\sum_{t=1}^T \mathrm{Var}(X_t)}.
\end{equation}

Effective sample size:

\begin{equation}
T_{\mathrm{eff}}^{(X)} = \frac{T}{\kappa_X}.
\end{equation}

\subsection*{8.2 Residual-Based Effective Sample Size}

For standardized residuals

\begin{equation}
r_t = \frac{y_t - \mu_t}{\sigma_t},
\end{equation}

\begin{equation}
\kappa_r
=
\frac{\mathrm{Var}\!\left(\sum r_t\right)}
{\sum \mathrm{Var}(r_t)}.
\end{equation}

\subsection*{8.3 Score-Based Effective Sample Size}

Define

\begin{equation}
\mathbf J_{\mathrm{diag}}
=
\sum_{t=1}^T \mathrm{Cov}(s_t).
\end{equation}

Define the Correlation Distortion Matrix (CDM):

\begin{equation}
\mathbf R
=
\mathbf J_{\mathrm{diag}}^{-1/2}
\mathbf J
\mathbf J_{\mathrm{diag}}^{-1/2}.
\end{equation}

If there is no temporal dependence in the score process,
$\mathbf R = \mathbf I$.

Scalar effective sample size measures can be derived from:

\begin{align}
\kappa_{\mathrm{tr}} &= \frac{1}{d}\mathrm{tr}(\mathbf R), \\
\kappa_{\mathrm{det}} &= \exp\left(\frac{1}{d}\log\det(\mathbf R)\right).
\end{align}

\section*{9. Decomposition of Information Distortion}

Define the Geometric Distortion Matrix:

\begin{equation}
\mathbf C_{\mathrm{geom}}
=
\mathbf H^{-1/2}
\mathbf J_{\mathrm{diag}}
\mathbf H^{-1/2}.
\end{equation}

Define the Correlation Distortion in Fisher coordinates:

\begin{equation}
\mathbf C_{\mathrm{corr}}
=
\mathbf C_{\mathrm{geom}}^{-1/2}
\mathbf C
\mathbf C_{\mathrm{geom}}^{-1/2}.
\end{equation}

Then:

\begin{equation}
\mathbf C
=
\mathbf C_{\mathrm{geom}}^{1/2}
\mathbf C_{\mathrm{corr}}
\mathbf C_{\mathrm{geom}}^{1/2}.
\end{equation}

Interpretation:

\begin{itemize}
\item $\mathbf C_{\mathrm{corr}}$ captures temporal correlation effects.
\item $\mathbf C_{\mathrm{geom}}$ captures mismatch between per-interval score variance and Gaussian Fisher curvature.
\item If $\mathbf C_{\mathrm{corr}} \approx \kappa \mathbf I$, distortion reduces to effective sample size scaling.
\item If $\mathbf C_{\mathrm{corr}}$ is anisotropic, geometric distortion remains even after correcting for $T_{\mathrm{eff}}$.
\end{itemize}

\section*{10. Relation to DCC and DIB}

If distortion is purely due to correlation:

\begin{equation}
\mathbf C \approx \kappa \mathbf I
\quad \Rightarrow \quad
\Sigma_{\mathrm{DCC}} \approx \kappa \mathbf H^{-1}.
\end{equation}

However, anisotropic distortion implies that
simple effective sample size corrections are insufficient.

Bias and variance corrections operate independently:

\begin{align}
\Sigma_{\mathrm{DCC}} &= \mathbf H^{-1} \mathbf J \mathbf H^{-1}, \\
\mathbf b_{\mathrm{DIB}} &= \mathbf H^{-1} \mathbf g.
\end{align}

\section*{11. Summary}

\begin{align}
\mathbf H &\rightarrow \text{Gaussian Fisher curvature} \\
\mathbf J &\rightarrow \text{Score fluctuation geometry} \\
\mathbf C &\rightarrow \text{Information distortion} \\
\mathbf R &\rightarrow \text{Correlation distortion} \\
\Sigma_{\mathrm{DCC}} &\rightarrow \text{Distortion-Corrected Covariance} \\
\mathbf b_{\mathrm{DIB}} &\rightarrow \text{Distortion-Induced Bias}.
\end{align}

Together these quantities distinguish temporal dependence,
geometric mismatch, variance inflation, and systematic bias.

\section*{8. Effective Sample Size and Temporal Dependence}

The Information Distortion Matrix $\mathbf C$ captures the total deviation between
Gaussian curvature and score fluctuations. However, this deviation may arise from two conceptually distinct mechanisms:

\begin{enumerate}
\item Residual temporal dependence between incremental contributions.
\item Geometric mismatch between instantaneous score variance and Gaussian Fisher curvature.
\end{enumerate}

We therefore analyze the variance-of-sum versus sum-of-variances decomposition.

\subsection*{8.1 Scalar Case: Variance Inflation}

For a scalar sequence $X_t$, define

\begin{equation}
L_X = \sum_{t=1}^T X_t.
\end{equation}

Then

\begin{equation}
\mathrm{Var}(L_X)
=
\sum_{t=1}^T \mathrm{Var}(X_t)
+
2\sum_{i<j}\mathrm{Cov}(X_i,X_j).
\end{equation}

Define the inflation factor

\begin{equation}
\kappa_X
=
\frac{\mathrm{Var}(L_X)}
{\sum_{t=1}^T \mathrm{Var}(X_t)}.
\end{equation}

If $X_t$ are independent, $\kappa_X=1$.
The effective sample size is

\begin{equation}
T_{\mathrm{eff}}^{(X)} = \frac{T}{\kappa_X}.
\end{equation}

\subsection*{8.2 Vector Case: Correlation Distortion Matrix}

For the score increments $s_t \in \mathbb R^d$, define

\begin{equation}
\mathbf J_{\mathrm{diag}}
=
\sum_{t=1}^T \mathrm{Cov}(s_t),
\end{equation}

which excludes cross-interval covariances.

Define the Correlation Distortion Matrix

\begin{equation}
\mathbf R
=
\mathbf J_{\mathrm{diag}}^{-1/2}
\mathbf J
\mathbf J_{\mathrm{diag}}^{-1/2}.
\end{equation}

\textbf{Proposition 1.}
If $s_t$ are temporally uncorrelated, then $\mathbf R = \mathbf I$.

\textbf{Proof.}
If $\mathrm{Cov}(s_i,s_j)=0$ for $i\neq j$, then
$\mathbf J = \mathbf J_{\mathrm{diag}}$.
Substitution yields $\mathbf R = \mathbf I$.
\hfill$\square$

Scalar effective sample size measures may be derived from invariants of $\mathbf R$:

\begin{align}
\kappa_{\mathrm{tr}} &= \frac{1}{d}\mathrm{tr}(\mathbf R), \\
\kappa_{\mathrm{det}} &= \exp\!\left(\frac{1}{d}\log\det(\mathbf R)\right).
\end{align}

If $\mathbf R \approx \kappa \mathbf I$, temporal dependence acts isotropically.

\section*{9. Geometric Distortion}

Even if temporal dependence is negligible, the instantaneous score variance
may not match Gaussian Fisher curvature.

Define

\begin{equation}
\mathbf C_{\mathrm{geom}}
=
\mathbf H^{-1/2}
\mathbf J_{\mathrm{diag}}
\mathbf H^{-1/2}.
\end{equation}

\textbf{Proposition 2.}
If instantaneous score variance matches Gaussian curvature,
i.e. $\mathbf J_{\mathrm{diag}} = \mathbf H$,
then $\mathbf C_{\mathrm{geom}} = \mathbf I$.

Thus $\mathbf C_{\mathrm{geom}}$ measures purely geometric mismatch.

\section*{10. Decomposition of Information Distortion}

The full Information Distortion Matrix can be decomposed as

\begin{equation}
\mathbf C
=
\mathbf C_{\mathrm{geom}}^{1/2}
\mathbf C_{\mathrm{corr}}
\mathbf C_{\mathrm{geom}}^{1/2},
\end{equation}

where

\begin{equation}
\mathbf C_{\mathrm{corr}}
=
\mathbf C_{\mathrm{geom}}^{-1/2}
\mathbf C
\mathbf C_{\mathrm{geom}}^{-1/2}.
\end{equation}

\textbf{Proposition 3.}
If distortion is purely due to temporal dependence and acts isotropically,
then

\begin{equation}
\mathbf C \approx \kappa \mathbf I,
\end{equation}

and

\begin{equation}
\Sigma_{\mathrm{DCC}} \approx \kappa \mathbf H^{-1}.
\end{equation}

\textbf{Interpretation.}
In this regime, a simple effective sample size correction
captures most of the distortion.

However, if $\mathbf C_{\mathrm{corr}}$ is anisotropic
or $\mathbf C_{\mathrm{geom}} \neq \mathbf I$,
then distortion cannot be reduced to scalar sample-size inflation.

\section*{11. Relation to Corrected Covariance and Bias}

Variance correction and bias correction arise from distinct mechanisms:

\begin{align}
\Sigma_{\mathrm{DCC}} &= \mathbf H^{-1} \mathbf J \mathbf H^{-1}, \\
\mathbf b_{\mathrm{DIB}} &= \mathbf H^{-1} \mathbf g.
\end{align}

Temporal dependence modifies $\mathbf J$.
Geometric mismatch modifies $\mathbf J_{\mathrm{diag}}$.
Score drift modifies $\mathbf g$.

Together, these quantities provide a complete second-order
description of algorithm-induced distortion.

\end{document}
